\documentclass[a5paper,10pt]{article}

% https://github.com/InDevRus/filippov-solutions
% Нумерация страниц
\pagenumbering{gobble}

% Настройка полей
\usepackage{geometry}
\geometry{left=1cm}
\geometry{right=1cm}
\geometry{top=1cm}
\geometry{bottom=1.5cm}

% Вставка таблицы
\usepackage{array}
\usepackage{tabularx}

% Инструменты выравнивания формул
\usepackage{amsmath}
\usepackage{mathtools}

% Диакритика в математике
\usepackage{amsfonts}

% Вычисления размеров на ходу
\usepackage{calc}

% Удобные записи производных
\usepackage[italic=true]{derivative}

% Настройки сносок
\usepackage{enotez}

% Длинные знаки векторов
\usepackage{anyfontsize}
\usepackage[b]{esvect}

% Вставка картинок
\usepackage{graphicx}

% Вставка красной строки
\usepackage{indentfirst}
\setlength{\parindent}{1em}

% Условия в командах
\usepackage{xifthen}

% Луа-скрипты
\usepackage{luacode}

% Настройка команд отдельно в математическом и текстовом режимах
\usepackage{mathcommand}

% Для повторения LaTeX кода
\usepackage{multido}

% Конфигурация отступов формул
\delimitershortfall-1sp
\usepackage{mleftright}
\mleftright

% Растягивание объектов
\usepackage{relsize}

% Обтекание текстом
\usepackage{wrapfig}
\usepackage{subcaption}
\usepackage{floatflt}

% Поддержка ввода кириллицы
\usepackage[english, russian]{babel}

% Настройка корневой позиции для компилятора
\usepackage{import}
\providecommand*{\ProjectRootPath}{..}

\import{\ProjectRootPath/preambles/macros/}{theorems}

\import{\ProjectRootPath/preambles/fonts}{font_setup}

% Знак градуса
\usepackage{siunitx}

\import{\ProjectRootPath/preambles/macros/}{operators}
\import{\ProjectRootPath/preambles/macros/}{redefinitions}

\import{\ProjectRootPath/preambles/fonts}{letters_setup}
\import{\ProjectRootPath/preambles/fonts/adjustments}{custom_superscripts}

\import{\ProjectRootPath/preambles/custom}{formatting}
\import{\ProjectRootPath/preambles/custom}{frames}
\import{\ProjectRootPath/preambles/custom}{list_setup}

% TODO: перевести команды на современный стиль объявления
\input{../preambles/plotting}

\begin{document}

$\der{y}{x} = \dfrac {\dot{y}} {\dot{x}} = \dfrac {2y} {x}$.
$\dfrac {\pow{x}{2} \primesub{y}{x} \br{x} - 2x y\br{x}} {\pow{x}{4}} = 0$.
$\vder{}{x} \br{\dfrac {y\br{x}} {\pow{x}{2}}} = 0$.
$y = C\pow{x}{2}$. Итак, траекториями системы являются параболы, 
проходящие через начало координат фазовой плоскости.

\begin{floatingfigure}{0.275\textwidth}
    \begin{tikzpicture}
        \pgfplotsset{
        VectorGraph/.style = {
                -stealth,
                quiver = {
                    u = x / (x^2 + y^2) ^ 0.5, 
                    v = 2 * y / (x^2 + y^2) ^ 0.5, 
                    scale arrows = 0.2}}
        }
        \begin{axis}[
            width = \textwidth,
            height = 7cm,
            ymin = -4.2,
            ymax = 4.7,
            xmin = -2.2,
            xmax = 2.45, 
            xtick = {-3, ..., 3},
            ytick = {-5, ..., 5},
            minor tick num = 3,
            declare function = {
                f(\C, \x) = \C * x^2;
            }]
            
            \addplot[domain=-3:3] {f(0.5, x)};
            \addplot[VectorGraph, samples = 6, domain = -2 : -0.5] {f(0.5, x)};
            \addplot[VectorGraph, samples = 6, domain = 0.5 : 2] {f(0.5, x)};
            \addplot[domain=-3:3] {f(-0.5, x)};
            \addplot[VectorGraph, samples = 6, domain = -2 : -0.5] {f(-0.5, x)};
            \addplot[VectorGraph, samples = 6, domain = 0.5 : 2] {f(-0.5, x)};
            \addplot[domain=-3:3] {f(1, x)};
            \addplot[VectorGraph, samples = 6, domain = -2 : -0.75] {f(1, x)};
            \addplot[VectorGraph, samples = 6, domain = 0.75 : 2] {f(1, x)};
            \addplot[domain=-3:3] {f(-1, x)};
            \addplot[VectorGraph, samples = 6, domain = -2 : -0.75] {f(-1, x)};
            \addplot[VectorGraph, samples = 6, domain = 0.75 : 2] {f(-1, x)};            
            \addplot[domain=-3:3] {f(2, x)};
            \addplot[VectorGraph, samples = 6, domain = -2 : -0.75] {f(2, x)};
            \addplot[VectorGraph, samples = 6, domain = 0.75 : 2] {f(2, x)};
            \addplot[domain=-3:3] {f(-2, x)};
            \addplot[VectorGraph, samples = 6, domain = -1.4 : -0.75] {f(-2, x)};
            \addplot[VectorGraph, samples = 6, domain = 0.75 : 1.4] {f(-2, x)};
            
        \end{axis}
    \end{tikzpicture}
\end{floatingfigure}

На рисунке показаны траектории вместе с направляющими векторами, посчитанными в некоторых 
точках вдоль траекторий. Характер направлений указывает на то, что решения удаляются от 
начала координат, то есть от нулевого решения $\Phi\br{t} = \begin{pmatrix} 0 \\ 0 \end{pmatrix}$,
так что решение неустойчиво, а особая точка 
$\tsys{x = 0 \\ y = 0}$ относится к типу <<узел>>.

Докажем по определению неустойчивость. Общее решение имеет вид 
$X\br{t} = \begin{pmatrix} A \pow{e}{t} \\ B \pow{e}{2t} \end{pmatrix}$.
Для $t \ge \sub{t}{0} = 0$
\begin{align*}
    \br{X\br{t} - \Phi\br{t}}^2 = \pow{A}{2} \pow{e}{2t} + \pow{B}{2} \pow{e}{4t} &&
    \br{X\br{\sub{t}{0}} - \Phi\br{\sub{t}{0}}}^2 = \pow{A}{2} + \pow{B}{2}
\end{align*}
\vspace{1mm}

Для $\epsilon = 1$ и любого $\delta > 0$ потребуем, чтобы $A$ и $B$ были таковы, 
что $0 < \pow{A}{2} + \pow{B}{2} < \pow{\delta}{2}$. Обе функции 
$\pow{e}{2t}$ и $\pow{e}{4t}$ возрастают, то есть при $t \ge \sub{t}{0} = 0$ будет
$\pow{e}{4t} \ge \pow{e}{2t}$. Однако, при
$t \ge \max\br{ \dfrac {1} {2} \ln\br{\dfrac {1}{\pow{A}{2} + \pow{B}{2}}}; 0}$
имеем
$\pow{A}{2} \pow{e}{2t} + \pow{B}{2} \pow{e}{4t} 
    \ge \pow{e}{2t} \br{\pow{A}{2} + \pow{B}{2}} 
    \ge 1 
    = \pow{\epsilon}{2}$,
и поэтому \linebreak решение неустойчиво.

\end{document}
